\section{担保 (Guaranteeing)}\label{sec:guaranteeing}

担保工作包涉及创建和分发相应的\emph{工作报告}，这需要满足某些条件。除了报告外，还需要一个证明验证者对其正确性承诺的签名。有了两个担保人签名，工作报告就可以分发给即将到来的 \Jam 链区块作者，以便在 $\xtguarantees$ 中使用，从而为担保人带来奖励。

我们假设在一个公开系统中，如果验证者出现故障并承诺了一份未能忠实代表对工作包应用 $\computereport$ 结果的报告，将受到严厉惩罚。总体而言，该过程如下：

\begin{enumerate}
    \item 评估工作包的授权，并与最近 \Jam 链状态中的授权池进行交叉引用。
    \item 评估每个工作项的 refine 逻辑。
    \item 根据纠删编解码器对工作包捆绑包和 refine 逻辑导出的数据进行分块。
    \item 组装和发布工作报告。
    \item 将上述分块分发到验证者集合中。
    \item 根据请求向其他验证者提供工作包捆绑包和导出数据，也有助于优化网络性能。
\end{enumerate}

对于担保人收到的任何工作包 $\mathbf{p}$，工作报告 $\mathbf{r}$（如果有）可通过以下方式确定：
\begin{equation}
  \mathbf{r} = \computereport(\mathbf{p}, c, \mathbf{l}, v)
\end{equation}

当需要 \Jam 链状态时，应始终使用最近区块的链状态。担保人可以选择 $c$（核心索引）、$\mathbf{l}$（segment-root 字典）和 $v$（保证验证者集合的大小）的值，但应注意选择允许工作报告在链上被包含的值（参见第 \ref{sec:workreportguarantees} 节）。通常 $c$ 应为当前分配给担保人的核心，$v$ 应为活跃验证者集合的大小，$\mathbf{l}$ 应从先前已担保的工作报告中导出。

在生成工作报告 $\mathbf{r}$ 之后，担保人可以安全地创建和分发有效载荷 $\tup{s, i}$，其中 $i$ 通过索引标识担保人，$s$ 是根据公式 \ref{eq:guarantorsig} 创建的签名。具体来说，$s$ 是使用验证者注册的 Ed25519 密钥对有效载荷 $l$ 的签名：
\begin{equation}
  l = \Xguarantee \concat \blake{\encode{\mathbf{r}}}
\end{equation}

为了最大化利润，担保人应要求工作报告满足第 \ref{sec:workreportguarantees} 节中描述的担保外在中生效的所有预期。这包括上下文有效性和授权池中授权的包含。不这样做不会导致惩罚，但会阻止区块作者包含该报告，从而减少奖励。

高级节点可以通过尝试预测报告到达区块作者时的链状态来最大化其报告可被链上包含的可能性。简单节点可在验证工作包时简单使用当前链头。为了最小化所做的工作，节点应在评估 $\Psi_R$ 函数以计算工作报告的工作摘要\emph{之前}进行所有此类评估。

在评估为值得担保的合理工作包后，担保人应通过尝试在核心上形成共识来最大化其工作不被浪费的机会。为此，他们应将工作包发送给同一核心上他们认为尚不知道该工作包的任何其他担保人。

为了最小化区块作者的工作量从而最大化预期利润，担保人应尝试从工作报告和包含自己以及尽可能多的其他人的证明集合中构建其核心的下一个担保外在。

为了最小化任何区块作者因反垃圾邮件措施而忽视担保人的机会，担保人每个时隙平均签名的工作报告不应超过两份。
